% Contextualizar metodos aprendidos em sala de converter sequencias infinitas em formulas
% Falar sobre a sequencia de fibonnaci e como definimos ela usando metodos anteriores
% Explicar o que e a diagonalizacao e Apontar como a diagonalizacao ẽ o mesmo que encontrar a funcao geradora nesse caso


Sequências numéricas são listas finitas ou infinitas de números que seguem ou não alguma regra ou padrão. Podemos defini-las de múltiplas maneiras, sendo as mais usadas em sala de aula as funções geradoras, o operador \textit{Cons} ($<$) e a definição de uma regra recursiva para a sequência.

A sequência de Fibonacci é uma das mais famosas da matemática, por ser usada em diversas áreas, como a biologia (no número de pétalas em flores, por tamanho) e a ciência da computação (no algoritmo de busca de Fibonacci e na \textit{heap} de Fibonacci), e por ter uma regra recursiva fácil de decorar: $F_{n+2}=F_{n+1}+F_n$, com $F_1=1$ e $F_0=0$. Também podemos defini-la por meio da função geradora $\frac{x}{1-(x+x^2)}$ ou, usando \textit{Cons}, por $F=0<(F+1<F)$.

A diagonalização é o processo em que convertemos uma matriz quadrada em um sistema composto por matrizes, de forma que $A=PDP^{-1}$, e os vetores ou matrizes multiplicados por esse sistema são redimensionados ao longo de direções específicas. A matriz original pode ou não representar um sistema de equações lineares multivariáveis. Neste trabalho, queremos mostrar as semelhanças entre a diagonalização e a função geradora da sequência de Fibonacci.
